\chapter{Projection, Interpretation, and Living Dynamics}
\label{ch:ch20-projection-interpretation}

This closing chapter asks a question that has been implicit since
Chapter~2 but never stated directly: what is lost whenever an object is
replaced by a description of it? Coordinates, projections, and Fourier
decompositions have each, in their turn, shown that a single object can be
described in more than one way, and that different descriptions make
different features easy to see. What none of the earlier chapters asked
explicitly is what a description \emph{cannot} see --- and it is that
question, posed with the ordinary linear algebra already developed in this
book, that occupies the pages that follow. The illustrative material in
this chapter is drawn from a debate in neuroscience about what it means
for a measured neural signal to "represent" something, but the
neuroscience is only a worked example. The mathematics being illustrated
is entirely general.

\section{Representation as a Mathematical Relation}

Recall from Chapter~2 that a vector's coordinates depend on a chosen
basis: the same displacement is recorded by different ordered pairs
under different coordinate frames, yet the vector itself, as a
displacement, does not change when the frame does. This distinction
generalizes. Let $X$ be a space of possible states of some system, and let
\[
\pi: X \to Y
\]
be any measurement or representation, associating to each state $x \in X$
a description $\pi(x)$ in some representational space $Y$. The description
$\pi(x)$ is not the state $x$ itself; it is $x$ viewed through the
particular relation $\pi$. Chapters~2, 13, and 14 each supplied an
instance of this same pattern --- coordinates, a projection onto a
direction, and a change of coordinate system --- without stating the
general principle they shared. This chapter states it, and asks what
follows from it.

\section{Projection and the Kernel}

Suppose $\pi$ is linear, as every projection encountered in this book has
been.

\begin{definition}
The \emph{kernel} of a linear map $P: X \to Y$ is $\ker(P) = \{x \in X :
P(x) = 0\}$.
\end{definition}

\begin{theorem}[Principle of Invisible Structure]
If $k \in \ker(P)$, then $P(x) = P(x+k)$ for every $x \in X$. Consequently,
no observation made solely through $P$ can distinguish $x$ from $x+k$.
\end{theorem}

\begin{proof}
By linearity, $P(x+k) = P(x) + P(k) = P(x) + 0 = P(x)$.
\end{proof}

The proof is one line, and deliberately so: the content of the theorem is
not its difficulty but its interpretation. The kernel of a projection is
not an absence of structure in the original system; it is the set of
directions in which the system can vary without the chosen measurement
noticing. Consider a biological illustration. Let a neuron's full state
be some high-dimensional $x$, encompassing its spike timing, its
neuromodulatory context, its metabolic condition, and its recent history,
and suppose an experiment measures only a single projection of this
state, a firing rate $P(x)$. Neuromodulatory context, metabolic
condition, and history may then inhabit, in whole or in part, the kernel
of $P$. Nothing about the neuron has disappeared; the measuring apparatus
simply cannot see those directions. The theorem gives this familiar
methodological caution --- absence of evidence is not evidence of absence
--- a precise mathematical location: the evidence is absent because it
lies in $\ker(P)$, and it can be found only by choosing a different $P$.

\section{Multiple Interpretability}

A second consequence follows from the same observation, stated now for the
image of $\pi$ rather than its kernel.

\begin{definition}
Given a measured value $y = \pi(x)$, an observer may define a
\emph{decoding map} $d: Y \to Q$ and interpret $d(y)$ as the quantity
represented by $x$. \emph{Multiple interpretability} is the fact that,
absent further constraint, many decoding maps $d_1, d_2, d_3, \dots$ may
all be defined on the same $Y$, so that a single measured state supports
several formally valid, but mutually incompatible, interpretations.
\end{definition}

If a measured value happens to equal some quantity $n$, it equally well
satisfies $y = 2(n+1) - 2$, or countless other relations; the raw
correlation between $y$ and $n$ does not, by itself, single out $n$ as
\emph{the} thing $y$ represents rather than any of the other quantities
$y$ is equally correlated with. This is worth distinguishing from a
related and more familiar idea: \emph{multiple realizability} asks how
many different physical systems can implement one formal structure.
Multiple interpretability asks the reverse question, how many formal
structures can be imposed on one physical system, and it is this second
question that bears on what a measurement is entitled to claim.

\section{Sparse Projections and Apparent Representation}

Modern data analysis often seeks a low-dimensional projection $z = Px$ of
a high-dimensional state $x \in \mathbb{R}^N$, with $P: \mathbb{R}^N \to
\mathbb{R}^k$ for $k \ll N$, sometimes further expressed as $x \approx
D\alpha$ for a sparse coefficient vector $\alpha$. Such projections can
genuinely reveal invariants, low-dimensional constraints, recurring modes,
and predictive relationships; none of that value is in question. What
does not follow from the existence of a successful sparse projection is
that the system being measured internally computes that projection, that
its coefficients are symbols manipulated by the system, or that the axes
the projection happens to use are privileged by the system itself rather
than convenient to the observer.

\begin{theorem}
Decodability does not imply internal representation.
\end{theorem}

This is not a theorem in the sense of the rest of this book, with a
derivation from prior results; it is closer to a definition made precise.
A decoder demonstrates that an observer, given $P(x)$, can recover some
target quantity. It does not, by that fact alone, demonstrate that the
system whose state is $x$ performs the same decoding, or uses $P(x)$ for
anything at all. The gap between these two claims is exactly the gap
between an observer's successful use of a projection and a claim about
what lies outside that projection's kernel.

\begin{sidebar}{A Neuroscientific Instance: Language and Reasoning}
A 2025 study of logical reasoning offers an unusually direct empirical
case of the gap this section describes. Kean and colleagues
\cite{kean2025} asked whether formal reasoning is carried out using the
brain's ordinary natural-language machinery, or whether language instead serves only as
an interface through which a problem is presented and a conclusion is
reported. Using functional imaging, they found that language-selective
brain regions were not substantially engaged while healthy participants
solved inductive and deductive reasoning tasks; a separate, domain-general
network associated with working memory and cognitive control was engaged
instead. They then studied two individuals with profound aphasia, whose
grammatical comprehension was severely impaired, and found that both
solved logical induction and matrix-reasoning problems at or above
typical levels despite this impairment. In the language of this chapter,
the natural-language system functions here as one measurable projection
of cognitive activity, $P_{\text{language}}$, correlated with reasoning
performance under ordinary conditions but not identical to whatever
operator actually carries out the reasoning; when that projection is
severely degraded, as in profound aphasia, the underlying capacity it had
seemed to track remains intact. A close correlation between a measured
quantity and an ability, in other words, is not by itself evidence that
the measured quantity is the mechanism producing that ability, whether
the system under study is a projected vector, a decoded neural signal, or
the relationship between language and thought.
\end{sidebar}

\section{Invariance Without Symbolism}

Chapter~1 showed that certain ratios are invariant under scaling; Chapter~6
showed that sine is odd and cosine is even under $\theta \to -\theta$. In
each case, invariance was a structural fact: a quantity $I(x)$ satisfying
$I(Tx) = I(x)$ for some transformation $T$. It is worth being precise
about what such a fact does and does not establish. If some measured
quantity $F$ satisfies $F(R(\theta)x) = F(x)$ for every rotation
$R(\theta)$ of Chapter~8, this establishes that $F$ is insensitive to
rotation --- a genuine structural discovery. It does not, by itself,
establish that $F$ "means" or "represents" the identity of the object
being rotated; that further claim requires more than invariance alone
provides. A productive habit, when this distinction threatens to
collapse, is to replace the question "what does this state represent"
with the question this book's tools can actually answer: which
transformations leave the measured relation unchanged?

\section{Fourier Decomposition and Observer Choice}

Chapter~19 showed that a periodic function may be written as a sum of
harmonics, $f(x) = a_0 + \sum_{n=1}^\infty (a_n\cos nx + b_n\sin nx)$.
Nothing in that construction requires the decomposition to be unique in
kind: the same function could equally be expressed relative to a basis of
localized pulses, wavelets, or modes derived directly from data, and each
choice of basis makes some features of the function easy to see and
others hard to see, exactly as Chapter~14 showed for polar versus
rectangular coordinates. That a signal admits a clean, useful Fourier
decomposition therefore does not imply that whatever physical system
produced the signal stores, computes, or manipulates a list of Fourier
coefficients; it implies only that the Fourier basis is a good basis for
that signal, in the same sense that polar coordinates are a good
coordinate system for a circle. The decomposition is a fact about the
relation between the signal and the chosen basis, not a fact about the
internal machinery of whatever produced the signal.

\section{The Commutativity Test}

The distinction between a useful decoder coordinate and a variable a
system actually uses suggests a test, rather than only a caution. Suppose
a measured state $x$ is projected to $z = P(x)$, and $z$ is claimed to be
a genuine variable used by the system rather than merely an observer's
coordinate. An intervention that changes $x$ to some $x'$ should then
produce a corresponding, predictable change from $z$ to $z' = P(x')$, and
this claim can be drawn as a square,
\[
\begin{array}{ccc}
x & \xrightarrow{\ P\ } & z \\[2pt]
\big\downarrow{\scriptstyle\text{intervention}} & & \big\downarrow{\scriptstyle ?} \\[2pt]
x' & \xrightarrow{\ P\ } & z'
\end{array}
\]
which commutes exactly when the effect of the intervention on $z$ agrees
with what $P$ predicts from its effect on $x$. If the square commutes
reliably under intervention, $z$ behaves as a variable genuinely coupled
to the system's dynamics. If it does not, $z$ may still be a perfectly
good coordinate for an observer to use, recoverable from the data and
correlated with behavior, without being a variable the system itself is
sensitive to. The test converts a hard interpretive question into an
ordinary question about whether a diagram commutes, which is exactly the
kind of question the algebra of this book is built to answer.

\section{The Representation Theorem}

The results of this chapter can be collected into a single closing
statement.

\begin{theorem}[The Representation Theorem]
Let $P: X \to Y$ be any non-injective representation of a system. Then
$\ker(P)$ is nontrivial, and $P(x) = P(x+k)$ for every $k \in \ker(P)$.
Consequently, no claim about distinctions lying entirely within
$\ker(P)$ can be established, supported, or refuted by observations made
through $P$ alone.
\end{theorem}

\begin{proof}
Non-injectivity means some two distinct states $x_1 \neq x_2$ satisfy
$P(x_1) = P(x_2)$; setting $k = x_1 - x_2$ gives a nonzero element of
$\ker(P)$ by linearity, so $\ker(P) \neq \{0\}$. The remaining claims are
the Principle of Invisible Structure of Section~20.2, restated for
general $k \in \ker(P)$.
\end{proof}

The theorem is almost embarrassingly simple to prove; its force lies
entirely in how often it is forgotten in practice, whenever the image of
a representation, $P(X)$, is treated as though it were $X$ itself.

\section{Further Directions}

The mathematics developed above --- kernels, invariance, basis-dependence,
and the commutativity test --- is enough to state the chapter's central
result rigorously. The same questions extend naturally beyond a single,
fixed projection. A living system does not simply sit still to be
measured: it acts on its environment and is in turn shaped by that
environment, so that organism and environment are better modeled as
co-evolving, $dO/dt = F(O,E)$ and $dE/dt = G(E,O)$, than as a fixed input
feeding a fixed output. More strikingly, the transformation rule
governing such a system need not itself stay fixed the way every operator
in Chapters~1 through 19 has: a system's own admissible transformations
can change over time, $T_{t+1} = \mathcal{U}(T_t, x_t, e_t)$, so that the
pair $(x_t, T_t)$, rather than $x_t$ alone, becomes the true state of the
system. Questions of viability (which states a system can occupy and
still persist), entrainment (how internal rhythms synchronize with
external ones without symbolic prediction), and information-theoretic
limits on what successive measurements can recover all extend from this
same starting point. Each is a substantial subject in its own right, and
each is a natural continuation of the mathematics in this chapter rather
than a further topic this book will develop; they are noted here as an
indication of where the kernel and commutativity ideas lead next, for a
reader who wants to follow them.

\section{Looking Back}

This book opened with two maps of the same city, one small enough for a
pocket and one large enough for a classroom wall, observing that a change
of scale alters measurement while leaving certain ratios untouched.
Chapter~1 called such ratios invariants. Every chapter since has, in its
own way, returned to the same pair of questions: what does a given
transformation preserve, and what does it discard? Coordinates showed
that one object admits many descriptions; rotation and the identities
built from it showed that different algebraic forms can encode identical
structure; projection showed that a measurement preserves some
distinctions while discarding others outright; polar coordinates and
amplitwists showed that the choice of representation changes what appears
simple; Euler's formula showed that separate formalisms can be
manifestations of a single underlying structure; Fourier decomposition
showed that even a decomposition into "elementary" pieces is relative to
a chosen basis. This chapter's contribution is the final piece of that
pattern, made precise rather than left as an intuition: every
representation is itself a transformation, and every transformation
carries a kernel, a set of distinctions it cannot see. Mathematics
supplies extraordinary power to describe structure. What this closing
chapter asks a reader to carry forward is the discipline to ask, of any
description powerful enough to seem complete, what its kernel contains.

\section{Exercises}
\begin{exercise}
Give two coordinate descriptions of the same vector using different
bases, and explain why the vector is not identical to either coordinate
list.
\end{exercise}

\begin{exercise}
Let $P: \mathbb{R}^2 \to \mathbb{R}$ be given by $P(x,y) = x$. Find
$\ker(P)$, and give two distinct vectors that $P$ cannot distinguish.
\end{exercise}

\begin{exercise}
A firing rate is computed by counting events over three different time
windows of different lengths. Explain why the resulting values need not
describe the same underlying feature of the system.
\end{exercise}

\begin{exercise}
Let $y = 2x+1$. Construct two different decoding maps from $y$ to a
target quantity, and explain why the correlation between $y$ and $x$
alone does not select one interpretation over the other.
\end{exercise}

\begin{exercise}
Let $P = \begin{pmatrix} 1 & 0 & 0 \\ 0 & 1 & 0 \end{pmatrix}$ project
$\mathbb{R}^3$ onto $\mathbb{R}^2$. Find $\ker(P)$, and explain why
successfully recovering a label from $P(x)$ does not prove that the
original system computes that label.
\end{exercise}

\begin{exercise}
Find a function of two variables that is invariant under the rotation
$R(\theta)$ of Chapter~8, and one that is invariant under the scaling
$S(s)$ of Chapter~15. Explain why exhibiting such an invariant is a
weaker claim than exhibiting a symbolic representation.
\end{exercise}

\begin{exercise}
Express a specific periodic signal using two different bases (for
instance, a Fourier basis and a basis of localized pulses), and discuss
whether the Fourier coefficients should be regarded as physically
separate components of the signal.
\end{exercise}

\begin{exercise}
Draw the commutativity-test square for a claimed decoder $z = P(x)$ of
your own construction, and describe what evidence would show the square
fails to commute under intervention.
\end{exercise}

\begin{exercise}
Explain, in your own words, why a useful decoder does not necessarily
reveal an internal code used by the system being decoded.
\end{exercise}

\begin{exercise}
Prove that an invertible change of coordinates has trivial kernel, and
explain why this means it preserves every distinction even though it
changes every coordinate description.
\end{exercise}

\begin{exercise}
Prove the Representation Theorem directly for a specific non-injective
linear map of your choice, exhibiting a nonzero element of its kernel.
\end{exercise}

\begin{exercise}[Capstone]
For a variable $z$ recovered from a measured system, distinguish three
separate claims: that $z$ is \emph{decodable} (an observer can recover
$z$ from the measurement), that $z$ is \emph{causally relevant} (altering
$z$ changes the system's subsequent behavior), and that $z$ is
\emph{represented} by the system (the system itself uses $z$ as a
variable in its own dynamics). Give an example in which the first claim
holds but the second does not, and an example in which the first two
claims hold but evidence for the third is still lacking. Explain, in a
short paragraph, why these three claims are routinely conflated, and why
separating them requires most of the mathematics developed in this book.
\end{exercise}
